\section{An aggregate replacement for singular root selection}
\label{sec:rootsets}

The Puiseux packet identifies a real obstruction, rather than a missing
estimate. For \(Y^2-X^2-a\), an exact Puiseux coefficient at the returned
ramification index distinguishes \(a=0\) from \(a>0\). Such a constructor
cannot be uniform in arbitrary fast coefficient names. The two existing
framework packets therefore construct simple branches with a supplied
derivative gap. Here we address a different output: the finite set of
roots, observed at a specified precision \cite{repo-puiseux,repo-contracts}.
Multiple roots and collisions
are allowed. We do not recover a singular Puiseux expansion.

The use of completed finite sets and multisets in constructive algebra
has an established antecedent in Richman's \emph{The fundamental theorem
of algebra: a constructive development without choice} (2000)
\cite{richman2000}. The
contribution here is the explicit first-order primitive recursive
construction, its finite precision budget, and its application to the
representation boundary recorded in this repository. In particular, the
set representation below deliberately forgets multiplicity; a theorem
about matching multisets would be a different result.

\subsection{The finite metric relation}

Write \(\mathbb Q[i]\) for the Gaussian rationals. For \(z=x+iy\), put
\[
 N(z)=x^2+y^2,\qquad \|z\|_1=|x|+|y|.
\]
The notation \(|z|\le r\), for rational \(r\ge0\), means the rational
comparison \(N(z)\le r^2\). Euclidean norms used in intermediate estimates
abbreviate their elementary squared-norm consequences; they are not
additional input objects. We use
\(N(zw)=N(z)N(w)\), \(|z|\le\|z\|_1\le2|z|\), and finite
triangle and Cauchy--Schwarz inequalities.

For nonempty finite lists \(S,T\) and rational \(r\ge0\), define
\[
 \operatorname{Cover}(S,T,r)
 \;:\Longleftrightarrow\;
 \forall z\in S\ \exists w\in T\; N(z-w)\le r^2.
\]
Put \(\operatorname{Close}(S,T,r)\) for the conjunction of both directed
covers. This is a bounded decidable predicate. Repeated list entries are
permitted and have no effect on the predicate. Thus this is precisely a
finite presentation of Hausdorff closeness of the underlying supports.

\begin{lemma}[Finite product stability]\label{lem:product-stability}
Let \(S=(z_1,\ldots,z_n)\) and \(T=(w_1,\ldots,w_n)\), where \(n\ge1\),
and suppose every entry has norm at most the rational \(R\ge1\). Write
\[
 F(X)=\prod_{j=1}^n(X-z_j),\qquad
 G(X)=\prod_{j=1}^n(X-w_j).
\]
Suppose corresponding coefficients below degree \(n\) differ in norm
by at most \(\delta\), where \(\delta\ge0\). Set
\(C=nR^{n-1}\). If \(C\delta\le r^n\), then
\(\operatorname{Close}(S,T,r)\).
\end{lemma}

\begin{proof}
For \(z\in S\), finite polynomial evaluation gives
\[
 |G(z)|=|G(z)-F(z)|
 \le\delta\sum_{i=0}^{n-1}|z|^i
 \le C\delta\le r^n.
\]
But \(G(z)=\prod_j(z-w_j)\). Test the finitely many rational comparisons
\(N(z-w_j)\le r^2\). If all fail, each squared distance is strictly
greater than \(r^2\), and multiplication yields
\(N(G(z))>r^{2n}\), a contradiction. Bounded search therefore produces
the required index. Exchange the lists for the reverse cover. The
argument also works when \(r=0\): a finite product of nonnegative
rationals can vanish only if a factor vanishes. Every search and
induction in this proof is on explicit finite data.
\end{proof}

No assignment pairing all roots injectively was used. For example,
\((0,0,1)\) and \((0,1,1)\) have identical supports, although their
multiplicities differ. The lemma must not be read as a multiplicity
matching theorem.

\subsection{A finite approximate-root dependency}

We reuse one explicit human-proof result from
\path{algebra/FTA_iPRA/FTA_iPRA.tex}, Section~4, ``Uniform descent
for a rational polynomial'' \cite{repo-fta}. For a nonconstant Gaussian-rational
polynomial of degree at most \(d\ge1\), its primitive recursive
\(\operatorname{Descent}\) operation satisfies
\[
 |P(z+\Delta)|\le q_d|P(z)|,
 \qquad q_d=1-\frac1{4D_d},\qquad D_d=2^d4^{2d^2}.
 \tag{RD}\label{eq:rational-descent}
\]
The supplied argument uses a bounded dyadic radius search, a dominant
term at three consecutive scales, and a bounded rational direction
calculation. It does not assume the existence of a root. This is the
only nontrivial algebraic-existence dependency of the construction below.
It has a written proof in the repository; we do not claim its complete
\texttt{Prf PRAAx} formalization is available.

Here is the precise finite corollary we require, with its proof. Given
\(0<\varepsilon\le1\), let
\[
 B=\max(1,\|P(0)\|_1),\quad
 K=\lceil B/\varepsilon\rceil,\quad J=4D_dK.
\]
Start at zero and perform exactly \(J\) descent steps, assigning the
zero step when the residual is already zero. Finite Bernoulli arithmetic
gives \(q_d^{4D_d}\le1/2\). Hence the output \(z\) satisfies
\[
 |P(z)|\le Bq_d^J\le B2^{-K}\le\varepsilon,
\]
because \(2^K\ge K+1\ge B/\varepsilon\). This defines a single
primitive recursive operation on rational polynomial codes, uniformly in
the degree and requested precision. Only rational polynomials occur;
there is no evaluation of an arbitrary primitive recursive program code.

For a slower but simpler executable interface, the same corollary yields
a completely bounded grid search. If \(P\) is monic and its lower
coefficients have norm at most an integer \(A\ge1\), any point with
residual at most one has norm at most \(R=A+2\). Indeed, for
\(u=|z|>A+2\), the finite geometric identity gives
\[
 |P(z)|\ge u^d-A\frac{u^d-1}{u-1}>1.
\]
Set \(U=2R+1\) and
\(L_P=\max(1,\sum_{i=1}^d i\|p_i\|_1U^{i-1})\).
Take \(m=\lceil4L_P/\varepsilon\rceil\). Use the corollary at accuracy
\(\varepsilon/2\), and round its output to the square grid of spacing
\(1/m\) in \([-R,R]^2\). The distance moved is at most \(1/m\), and
the finite power-difference identity bounds the residual change by
\(L_P/m\le\varepsilon/4\). Thus the grid contains a point of residual
at most \(\varepsilon\). Scanning its \((2Rm+1)^2\) points is another
uniform PR approximate-root algorithm. This proof of grid completeness
depends on (RD), not on a classical compactness or FTA principle.

\subsection{Uniform finite factorization without nested coefficient names}

For a polynomial with Gaussian-rational coefficients, let
\(\|P\|_{\mathrm{coef},1}=\sum_i\|p_i\|_1\).
This norm is rational, subadditive, and submultiplicative, as follows
directly by expanding the finite convolution.

\begin{theorem}[Uniform rational approximate factorization]
\label{thm:approx-factor}
Let \(P\in\mathbb Q[i][X]\) be monic of degree \(n\ge1\), and let
\(0<\eta\le1\) be rational. From these finite inputs a primitive
recursive operation returns \(z_1,\ldots,z_n\in\mathbb Q[i]\) such that
\[
 \left\|P-\prod_{j=1}^n(X-z_j)\right\|_{\mathrm{coef},1}\le\eta.
\]
The operation is uniform in \(n\). Its mathematical proof uses (RD).
\end{theorem}

\begin{proof}
Choose an integer \(A_0\ge1\) bounding the \(\ell^1\) norm of every
coefficient of \(P\). Before performing any deflation, compute
\[
 A_{j+1}=(A_j+1)(2A_j+5)^n\quad(0\le j<n),\qquad
 L=\prod_{j=0}^{n-1}(2A_j+5),\qquad
 \varepsilon=\frac{\eta}{2nL}.
 \tag{AF}\label{eq:factor-budget}
\]
These are finite integer recursions. Let \(P_0=P\). At stage \(j<n\),
\(P_j\) is monic of degree \(n-j\). Apply the approximate-root
operation at tolerance \(\varepsilon\) and obtain \(z_{j+1}\). Perform
synthetic division, an exact finite recurrence, to obtain
\[
 P_j(X)=(X-z_{j+1})P_{j+1}(X)+r_j,
 \qquad r_j=P_j(z_{j+1}).
 \tag{SD}\label{eq:synthetic-division}
\]
The quotient is monic of degree \(n-j-1\).

Inductively, every coefficient of \(P_j\) has \(\ell^1\) norm at most
\(A_j\). The just-proved residual radius bound gives
\(|z_{j+1}|\le A_j+2\), because \(\varepsilon\le1\). Therefore
\(\|z_{j+1}\|_1\le2A_j+4\). The synthetic recurrence, beginning with
leading quotient coefficient one, bounds each quotient coefficient by
\[
 (A_j+1)(1+\|z_{j+1}\|_1)^n
 \le (A_j+1)(2A_j+5)^n=A_{j+1}.
\]
For clarity, the elementary recurrence estimate is obtained by repeatedly
expanding \(b_{k-1}=p_k+z_{j+1}b_k\): its right side is a sum of at
most \(n\) terms, each dominated by a term of
\((A_j+1)(1+\|z_{j+1}\|_1)^n\). Thus it does not require a loop
invariant about infinitely many coefficients.

After all \(n\) divisions, \(P_n=1\). Substituting
\eqref{eq:synthetic-division} repeatedly gives the exact identity
\[
 P(X)-\prod_{k=1}^n(X-z_k)
   =\sum_{j=0}^{n-1}r_j\prod_{k=1}^{j}(X-z_k).
 \tag{TE}\label{eq:factor-telescoping}
\]
Here the empty product is one. Since \(\|r_j\|_1\le2\varepsilon\),
submultiplicativity gives
\[
 \left\|P-\prod_{k=1}^n(X-z_k)\right\|_{\mathrm{coef},1}
 \le 2\varepsilon\sum_{j=0}^{n-1}
       \prod_{k=0}^{j-1}(2A_k+5)
 \le2\varepsilon nL=\eta.
\]
Each division manipulates a finite rational coefficient list; the outer
loop has exactly \(n\) stages. Consequently this construction avoids the
higher-type iteration issue in repeatedly deflating exact Cauchy names.
Its cost is large but primitive recursive.
\end{proof}

\subsection{A uniform finite-set name for all roots}

Fix one primitive recursive Gaussian-rational coefficient evaluator
\(a(i,s,\mathbf p)\). Parameters \(\mathbf p\), degree \(n\ge1\), and
precision indices are ordinary numerical inputs. Assume, for \(i<n\),
\[
 |a(i,s,\mathbf p)-a(i,t,\mathbf p)|\le2^{-s}+2^{-t}.
 \tag{CN}\label{eq:coefficient-name-root}
\]
The leading coefficient is exactly one. Suppress \(\mathbf p\) below.
Let
\[
 A=3+\left\lceil\max_{i<n}\|a(i,0)\|_1\right\rceil,
 \quad R=A+2,\quad C=nR^{n-1}.
\]
Then every sampled coefficient has Euclidean norm at most \(A-1\).
In the approximate-factorization procedure one can use \(A_0=2A\),
which bounds every sampled coefficient's \(\ell^1\) norm.
Choose an integer \(b\) with \(2^b\ge4C\), by bounded binary-length
calculation, and set
\[
 E(t)=n(t+2)+b,\qquad
 P_t(X)=X^n+\sum_{i<n}a(i,E(t))X^i.
\]
Apply Theorem~\ref{thm:approx-factor} to \(P_t\) with
\(\eta_t=2^{-E(t)}\). Denote its output list by \(S_t\) and its exact
product polynomial by \(F_t\).

\begin{theorem}[Collision-tolerant root-support names]
\label{thm:root-support}
The lists \(S_t\) form a primitive recursive finite-set name satisfying
\[
 u\ge t\quad\Longrightarrow\quad
 \operatorname{Close}(S_t,S_u,2^{-t-2}).
 \tag{RS}\label{eq:root-support-fast}
\]
All entries of all lists have norm at most \(R\). The construction is
uniform in the numerical degree and parameters of the fixed evaluator.
Its coefficient approximation and finite Hausdorff assertions have
\(\ipra\) proofs relative to the written descent lemma (RD).
\end{theorem}

\begin{proof}
Each lower coefficient of \(F_t\) has norm at most \(A\), because its
distance from the sampled coefficient is at most \(\eta_t\le1\).
Every entry of \(S_t\) is an exact rational root of the monic rational
polynomial \(F_t\). The radius argument above, now with residual zero,
therefore gives the common bound \(R=A+2\). This bound comes from the
\emph{final product coefficients}; we do not confuse it with the much
larger temporary bounds \(A_j+2\) used during deflation.

For \(u\ge t\), corresponding coefficients of \(F_t,F_u\) differ by
at most
\[
 2^{-E(t)}+(2^{-E(t)}+2^{-E(u)})+2^{-E(u)}
 \le4\cdot2^{-E(t)}.
\]
Consequently
\[
 C\,4\cdot2^{-E(t)}
 \le2^{-n(t+2)}=(2^{-t-2})^n.
\]
Lemma~\ref{lem:product-stability} proves \eqref{eq:root-support-fast}.
No equality or separation decision about represented coefficients or
roots occurs anywhere in the construction.
\end{proof}

The exact retained information is a nonempty finite-set name, together
with monic product polynomials converging coefficientwise to the input.
The next proposition proves that its zero-distance locus is exactly the
root locus, with an explicit finite-precision argument.

\begin{proposition}[Adequacy of the root-support representation]
\label{prop:root-adequacy}
Let \(z\) be a represented point, and supply an integer bound
\(Z\ge\max(R,1)\) with \(|z|\le Z\). Put
\[
 D_t(z)=\min_{w\in S_t}|z-w|,\quad
 C_Z=nZ^{n-1},\quad
 M_Z=A\sum_{i=1}^n iZ^{i-1}.
\]
The distance is a represented nonnegative real, computable by finite
minimum and square-root operations with prescribed precision. Define
\[
 \operatorname{Member}(z,S)\;:\Longleftrightarrow\;
 \forall k\ \exists t\ge k\;D_t(z)\le2^{-k}.
\]
Then \(\operatorname{Member}(z,S)\) holds if and only if \(P(z)=0\).
The forward production of indices from a root is primitive recursive.
Moreover the following error bounds hold:
\begin{align}
 |F_t(z)-P(z)|&\le2C_Z2^{-E(t)},\label{eq:root-eval-error}\\
 |P(z)|&\le M_ZD_t(z)+2C_R2^{-E(t)}.
 \label{eq:root-distance-residual}
\end{align}
\end{proposition}

\begin{proof}
Each lower coefficient of \(F_t\) differs from the represented input
coefficient by at most \(2\cdot2^{-E(t)}\): one contribution is the
approximate-factorization error and one is the sampled-name error.
The finite sum of powers at \(|z|\le Z\) proves
\eqref{eq:root-eval-error}. The power-difference identity proves
\(|P(z)-P(w)|\le M_Z|z-w|\) on this disk. At \(w\in S_t\),
\(F_t(w)=0\), and \(|w|\le R\), so
\(|P(w)|\le2C_R2^{-E(t)}\). This proves the inequality in
\eqref{eq:root-distance-residual} for each individual distance
\(|z-w|\), hence also for their finite minimum. No exact minimizing
index for represented distances has to be selected. For example a
finite minimum of real names can be implemented with
\(\min(u,v)=(u+v-|u-v|)/2\), with each arithmetic precision prescribed.

For completeness, here is a fully finite precision budget in the
root-to-membership direction. Let \(z(s)\) be its rational samples,
with sample error at most \(2^{-s}\) in represented notation. Set
\[
 U=Z+1,\quad C_U=nU^{n-1},\quad
 M_U=A\sum_{i=1}^n iU^{i-1},\quad T=C_U+M_U,
\]
and let \(\ell(v)\) be any primitive recursive integer choice with
\(2^{\ell(v)}\ge\max(1,v)\), for example binary length after rounding
\(v\) upwards. The rational polynomial evaluations
\[
 V_s=z(s)^n+\sum_{i<n}a(i,s)z(s)^i
\]
have error at most \(T2^{-s}\) relative to \(P(z)\); this follows by
separately varying the coefficients and the point. Thus \(P(z)=0\)
gives \(|V_s|\le T2^{-s}\). This is an explicit normalized
vanishing-error formulation of the root equation, not an invocation
of a limit principle. Direct coefficient comparison now gives the
entirely finite inequality
\[
 |F_t(z(s))|
 \le 2C_U2^{-E(t)}+(C_U+T)2^{-s}.
 \tag{FE}\label{eq:finite-root-evaluation}
\]
Given \(k\), choose
\[
 t=k+1+\ell(2C_U),\qquad
 s=\max\{k+1,\ n(k+3)+\ell(C_U+T)\}.
\]
The first summand in \eqref{eq:finite-root-evaluation} is at most
\(2^{-n(k+3)}\), because
\(E(t)\ge n(k+3)+\ell(2C_U)\); the second has the same bound.
Their sum is at most \(2^{-n(k+1)}\). All points in this last
calculation are rational. Testing their finitely many squared
distances and applying the product argument produces \(w\in S_t\)
with \(|z(s)-w|\le2^{-k-1}\). Since \(s\ge k+1\), adding the
sample error gives \(|z-w|\le2^{-k}\), as required. Both indices
and the successful finite search are fixed PR constructions. In a
pure sample formulation, the displayed represented inequalities are
obtained by substituting these precisions and the ordinary rational
slack-removal lemma; no real trichotomy is used.

Conversely, assume membership. Given a residual target \(2^{-k}\),
take \(j=k+\ell(M_Z+2C_R)\), and instantiate membership to obtain
\(t\ge j\) with \(D_t(z)\le2^{-j}\). Since \(E(t)\ge t\),
\eqref{eq:root-distance-residual} yields
\[
 |P(z)|\le(M_Z+2C_R)2^{-j}\le2^{-k}.
\]
This proves the root equation at every requested precision. The
existential index is used by intuitionistic existential elimination;
no unbounded search for it is assumed.
\end{proof}

By \eqref{eq:root-support-fast}, the distances \(D_t(z)\) have the
same explicit Cauchy error bound as the sets. Accordingly the displayed
membership formula is equivalent to zero distance from the represented
set. The generic finite-set chain construction supplies actual PR point
names when such a realization is needed. No arbitrary infinite choice
of labels is implicit in the root-support theorem.

\subsection{What the new representation resolves}

Lemma~\ref{lem:product-stability} is also a finite perturbation theorem.
For two monic input polynomials whose sampled coefficient differences
are at most \(\delta\), with splitting errors \(\eta,\eta'\), lists
in a common radius \(R\) are \(r\)-close whenever
\[
 nR^{n-1}(\delta+\eta+\eta')\le r^n.
 \tag{CP}\label{eq:coefficient-perturbation}
\]
This gives an explicit H\"older precision budget even at repeated roots.
For coefficients depending on another represented variable, combine a
supplied coefficient continuity modulus with
\eqref{eq:coefficient-perturbation}; this produces a modulus for the
root-support map. Neither a simple-root derivative bound nor root
separation is required.

For example, the obstruction family \(Y^2-X^2-a\) has a uniformly
represented root support at every represented input pair \((X,a)\).
For real \(X\), \(a\ge0\), the familiar semantic description is
\(\{-\sqrt{X^2+a},\sqrt{X^2+a}\}\); at \((0,0)\) the support is the
singleton \(\{0\}\). Coefficient perturbations of order \(\varepsilon^2\)
suffice for root-support accuracy of order \(\varepsilon\). The theorem
never has to decide whether the two roots coincide. It therefore repairs
the output contract for tasks needing distances, images, approximate
optimization, or coverage of all roots. It does not provide individually
accessible singular Taylor/Puiseux coefficients, multiplicity matching,
or a continuously labeled root on a parameter domain.

\subsection{Formalization obligations}

The file \path{lean/RootSetTargets.lean} defines Gaussian-rational
arithmetic, finite products, coefficient error bounds, the bounded cover
predicate, and the complete finite approximate-root constructor from a
supplied descent map with the displayed \(q_d,J\) budget. It also defines
point-name evaluation, normalized residual equality, and membership by
rational sample coverage. The forward adequacy indices are actual
definitions. Six named propositions specify the remaining correctness
proofs, including the descent-to-root bridge and the adequacy equivalence.
These propositions are \emph{unproved targets}, not axioms. The intended order
is: (1) polynomial convolution and synthetic-division identities;
(2) finite squared-norm inequalities and product stability;
(3) the explicit (RD) proof from the FTA packet;
(4) bounded approximate-root production; (5) the quantifier-free
deflation invariant and telescoping error bound; and (6) the fixed-name
precision substitution proving \eqref{eq:root-support-fast}.

For an object-theory proof, encode lists and Gaussian rationals as natural
numbers. Replace every list quantifier by its finite Boolean fold, retain
\eqref{eq:coefficient-name-root} as a hypothesis instantiated only at
the finitely needed indices, and induct on the finite deflation table.
The invariant includes coefficient bounds, degree, monicity, the exact
partial telescoping identity, and the accumulated rational error bound.
An ordinary Lean theorem over \texttt{Rat} is a useful preliminary layer;
the final goal is a derivation in the embedded \(\ipra\) calculus.
The reference program checks finite certificates with exact fractions.
It does not purport to execute the enormous worst-case descent schedule
or to verify the complete object-theory derivation.
